% Input fragment for the exact finite-algebra specification's bounds.
\subsection{Explicit degree and coefficient bounds for the specification}
\label{subsec:finite-algebra-bounds}

The following bounds concern the expanded polynomial defined by the
indexed determinant formula. They do not require expanding its matrices
or its final finite product. For a polynomial $f$, let $\|f\|_1$ denote
the sum of the absolute values of its coefficients. Then
$\|fg\|_1\le\|f\|_1\|g\|_1$ and
$\|f+g\|_1\le\|f\|_1+\|g\|_1$.

Give every base coordinate weight one, and assign a positive integer
weight $w_i$ to each algebra generator $z_i$. Choose the weights so that
every nonleading term of its monic relation has weight at most
$e_iw_i$. Reduction by these relations never increases weighted degree.
Writing $W=\sum_i(e_i-1)w_i$, the entries of a multiplication matrix
$M_f$ therefore have base degree at most $\operatorname{wdeg}(f)+W$:
the input basis monomial has weight at most $W$, and the output basis
monomial has nonnegative weight.

For a matrix over the polynomial ring, use the maximum column sum of
the coefficient norms of its entries. This norm is submultiplicative.
If
\[
 z_i^{e_i}+\sum_{k<e_i}c_{ik}z_i^k=0,
\]
the block companion description gives
\[
 \|M_{z_i}\|\le
 \max\left\{1,\sum_{k<e_i}\|M_{c_{ik}}\|\right\}.
\]
Extending to later monic generators preserves this estimate, since
multiplication by an earlier element is repeated on diagonal blocks,
up to a simultaneous permutation of rows and columns. Thus weights
and bounds $\|M_{z_i}\|\le2^{b_i}$ can be computed recursively from
the coefficient DAG. The source compiler records the entire table,
and the independent checker verifies the corresponding inequalities.

For the point quartics the weights can be taken to be two. A coefficient
bound for the companion matrix is the sum of the bounds for its two
nonzero lower coefficients. For each of the subsequent quadratic
generators, its nonzero branch has only the constant coefficient
$LM$, so one may take
\[
 w_i\ge\tfrac12\operatorname{wdeg}(LM),\qquad
 b_i\ge\log_2\|M_{LM}\|.
\]
The zero branch has relation $z_i=0$ and satisfies the same bounds.
Consequently one uniform table applies to all $2^{86}$ fixed branches.

Put $n=2^{508}$ for a uniform rank bound, $m=297$, and
$T=m n(n+1)/2$ for the largest separating parameter. Suppose the
recorded bounds give
\[
 \operatorname{wdeg}(G)\le w_G,\quad \|M_G\|\le2^{b_G},\qquad
 \operatorname{wdeg}(E_i)\le w_E,\quad \|M_{E_i}\|\le2^{b_E}.
\]
Set
\begin{align*}
 b_\lambda&=\lceil\log_2(m+1)\rceil
       +\max_i\{i\lceil\log_2 T\rceil+b_i\},\\
 \alpha&=w_G+\max_i w_i+W,\qquad \beta=b_G+b_\lambda,\\
 \alpha_E&=w_E+W.
\end{align*}
These bound the entry degrees and operator norm of
$M_{G\lambda_\ell}$, uniformly for all prescribed $\ell$ and all
source branches.

The coefficients of $H(V)=\det(I-VM)$ have base degree at most
$A_H=n\alpha$, and its full coefficient norm is at most
$2^{B_H}$ with $B_H=n(\beta+1)$. Indeed the coefficient norm of a
determinant is bounded by the product of its column sums. This bound
already includes the sum over determinant permutations; no additional
factorial is needed. By multilinearity in columns,
\[
 R_i(V)=\left.\frac{\partial}{\partial t}
       \det(I-VM+tM_{E_i})\right|_{t=0}
\]
has base coefficient degree at most
$A_R=(n-1)\alpha+\alpha_E$ and coefficient norm at most $2^{B_R}$,
where
\[
 B_R=\lceil\log_2 n\rceil+b_E+(n-1)(\beta+1).
\]
This accounts for the $n$ possible columns containing $M_{E_i}$.

For each degree case $d\le n$, the Sylvester matrix for
$\operatorname{Res}(H_d,H'_d)$ has $2d-1$ rows. Every row has
coefficient norm at most $d\|H\|_1$. Hence uniform bounds for the
reciprocal scale $L=a_d^3\operatorname{Res}(H_d,H'_d)$ are
\begin{align*}
 A_L&=(2n+2)A_H,\\
 B_L&=(2n+2)B_H+(2n-1)\lceil\log_2 n\rceil.
\end{align*}
The reciprocal residual has degree at most $n(A_L+1)$ and coefficient
norm at most $2^{nB_L+\lceil\log_2(n+1)\rceil}$. The reversed trace
residuals have degree at most $A_R+(n-1)(A_L+1)$ and coefficient norm
at most $2^{B_R+(n-1)B_L}$. Together with the high-coefficient
conditions $a_k=0$, these give a bound for the sum of squares by
doubling the largest degree, doubling the largest logarithmic norm
bound, and adding $\lceil\log_2(n+87)\rceil$ to the latter.
For the reciprocal estimate, no division of coefficient norms is used:
its $j$th coefficient is the polynomial
$a_{d-j}a_d^{3j-1}\operatorname{Res}(H_d,H'_d)^j$.
Its degree is at most $jA_L$ and its logarithmic norm bound is at most
$j(3B_H+B_{\rm Res})=jB_L$.

The number of final factors is also explicit. Write $n_0=4^{211}$.
For a mask with $k$ nonzero quadratic branches, its rank is
$n_k=n_0 2^k$. There are $\binom{86}{k}$ such masks, and its number
of separating-parameter and degree choices is
$n_k(1+m n_k(n_k+1)/2)$. Including the $212$ point charts and $256$
height choices gives exactly
\[
 B_{\mathrm{tot}}=54272\left[
 n_0 3^{86}+
 \frac{297}{2}\left(n_0^3 9^{86}+n_0^2 5^{86}\right)\right]
\]
factors. Multiplying the uniform branch degree and logarithmic
coefficient-norm bounds by $B_{\mathrm{tot}}$ bounds the complete
polynomial. All quantities here are exact integers, apart from the
displayed ceiling notation, which is evaluated by integer bit lengths.

For the emitted coefficient specification the resulting convenient
rounded bounds are
\[
 \deg P\le2^{3122},\qquad H(P)\le2^{\,2^{3127}}.
\]
The machine-readable specification records the stronger unrounded
integer degree bound and the stronger integer bound on $\log_2 H(P)$.
These numbers are rigorous upper bounds for the literal indexed
polynomial, not measurements of an expanded coefficient table.
